\documentclass[mnsc]{informs4}
\usepackage{eqndefns-left}
\RequirePackage{tgtermes}
\OneAndAHalfSpacedXII
\usepackage{amsmath,amssymb,graphicx,booktabs,tabularx}
\usepackage{natbib}
\bibpunct[, ]{(}{)}{,}{a}{}{,}%
\def\bibfont{\small}%
\def\bibsep{\smallskipamount}%
\def\bibhang{24pt}%
\def\newblock{\ }%
\def\BIBand{and}%
\usepackage{hyperref}
\hypersetup{colorlinks=true,linkcolor=black,citecolor=black,urlcolor=blue}
\graphicspath{{overleaf/figures/}{../figures/}}
\EquationsNumberedThrough
\TheoremsNumberedThrough
\ECRepeatTheorems
\RUNAUTHOR{Nguyen}
\RUNTITLE{Review Allocation and AI Adoption}
\TITLE{Who Reviews After AI Adoption? The Distribution and Continuity of Code Review Work}
\ARTICLEAUTHORS{%
\AUTHOR{Minh Hoang Nguyen}
\AFF{School of Business and Technology Management, KAIST,\
\EMAIL{ngmhoang2306@gmail.com}}
}
\MANUSCRIPTNO{MNSC-0001-2026.00}
\ABSTRACT{%
General-purpose technologies reshape work less by replacing people than by rearranging which people do which tasks. We study that rearrangement in code review, a complementary task an AI coding agent makes more necessary without performing. Theory on knowledge-based hierarchies predicts the direction: because matching problems to solvers is costly, routine problems are handled broadly and exceptional ones route to specialists, so a technology absorbing routine production should concentrate residual evaluation on a project's most knowledgeable reviewers. Testing this on public GitHub repositories around first detected AI adoption yields a methodological result that matters for this literature. Review concentration is a bounded, strongly mean-reverting share, so matched designs that balance repository characteristics but not the baseline outcome recover a substantial mechanical effect: a placebo using control repositories alone, where the true effect is zero by construction, reproduces roughly half of a naive estimate. Unit fixed effects do not remove it, and pre-trend tests do not detect it, because reversion acts on levels while those tests examine slopes. We re-identify the estimate on a design matched on the baseline outcome. Separately, within repositories and months, machine-authored pull requests attract more human reviewers and are markedly less likely to merge unreviewed than comparable human-authored ones, while dependency bots show the opposite pattern. Evaluation work is actively reallocated, not passively inherited.
}

% Keep the template's metadata panel blank until dates and publication details exist.
\makeatletter
\def\ABSfont{\EGT}
\def\theARTICLEABSTRACT{%
  \HOOKb
  \vspace*{18pt}
  \begin{minipage}[t]{.3\textwidth}\parindent1em\ABSfont\noindent\mbox{}\end{minipage}\hfill
  \begin{minipage}[t]{.65\textwidth}\parindent1em
    \ABSfont
    \noindent\theABSTRACT\endgraf
    \vskip5pt
    \theFUNDING
    \theKEYWORDS
    \theSUBJECTCLASS
    \theAREAOFREVIEW
    \theMSCCLASS
    \theORMSCLASS
    \if@BLINDREV\else\theHISTORY\fi
    \noindent\hrulefill
  \end{minipage}%
  \vspace*{0pt}
}
\makeatother
\KEYWORDS{Generative AI, Code Review, Work Allocation, Open-Source Software}
\begin{document}
\maketitle
\section{Introduction}

General-purpose technologies rarely change work by substituting machines for people in a single task. They change it by rearranging which people perform which tasks, and that rearrangement, rather than the direct substitution, is usually where the organizational consequences lie \citep{bresnahan2002information,acemoglu2019automation}. The canonical demonstrations are structural rather than aggregate. Computed-tomography scanners did not simply speed up imaging; they renegotiated the working relationship between radiologists and technologists, and the same machines produced divergent organizational forms in different departments \citep{barley1986technology}. Information technology in the workplace shifted the salient question from a division of labor to what \citet{zuboff1988smartmachine} termed a division of learning. Robotic surgery removed the practical need for a resident's hands, concentrating operative work on senior surgeons and sharply reducing the practice through which surgical skill had been acquired \citep{beane2019shadow}.

AI coding agents present the same class of question in a new setting. An agent can draft a pull request, but it does not evaluate, explain, revise, or accept one. Evaluation remains human, and a technology that raises the volume and changes the character of proposed changes alters the evaluation task without performing it. The organizational question is therefore not whether coding becomes faster. It is how the complementary evaluation work is redistributed across the people who remain responsible for it.

Theory supplies a direction rather than an open set of possibilities. In a knowledge-based hierarchy, agents specialize because matching a problem to someone who can solve it is costly: those nearer production learn the common and easy problems, while a smaller set of specialists absorbs the rare and hard ones \citep{garicano2000hierarchies}. A technology that absorbs routine production does not leave the residual problem distribution unchanged. It removes the easy cases disproportionately, so the evaluation work that remains is weighted toward exceptions, and exceptions route to the most knowledgeable evaluators. Concentration of review on a project's busiest reviewers is thus the prediction of an existing model, not a conjecture about AI in particular. \citet{barley1986technology} supplies the necessary qualification: the structural response to an identical technology depends on local social context, so the relevant object is a distribution of responses, not one average effect.

Code review is an unusually good setting for testing this. It is a genuine complementary task, performed by identifiable people, that detects defects and transfers knowledge about a changing code base \citep{bacchelli2013expectations}. It is also the mechanism by which projects distribute practice: reviewing is how participants other than the author learn what is in the system. Where the studies above infer structural change from ethnographic observation or from firm-level aggregates, code review leaves a complete record of individual acts: every review, by whom, on which change, at what time. The restructuring of a work process can be measured directly rather than inferred. Technologies that permit work to be distributed at the task level are known to generate new practices for routing and filtering it \citep{leonardi2008transformational}; here those practices are observable.

Measuring it directly, however, turns out to require more care than this literature has generally exercised, and establishing that is our first contribution. Review concentration is a bounded share that also mean-reverts strongly: repositories that are unusually concentrated in one period are less so in the next, whether or not anything was done to them. A matched design that balances repository characteristics but not the baseline level of the outcome therefore compares arms that would converge even under no treatment. We quantify the resulting bias with a placebo constructed from control repositories alone, assigning a fabricated treatment so that the true effect is zero by construction and calibrating the fabricated imbalance to the real one. Roughly half of a naive concentration estimate is reproduced this way. Unit fixed effects do not remove the problem, because they absorb differences in levels while reversion makes the expected change a function of the level; and the conventional pre-trend test does not detect it, because it examines pre-period slopes while the operative difference is in levels. This failure mode applies to any difference-in-differences estimate of a bounded, reverting share, rate, or index, and the diagnostic for it is straightforward once stated.

We therefore re-identify the concentration estimate under a design built for it: matching on the baseline outcome itself, within a population of repositories whose reviewer pool is large enough that the top-two share is not mechanically pinned at its upper bound, and over a fixed calendar window that permits controls to be paired with any treated repository. The placebo returns approximately zero when the baseline gap is closed, so a balanced design removes the mechanical channel and what remains is interpretable. [RE-MATCHED ESTIMATE PENDING: insert the re-estimated concentration effect, its interval, and the balanced-design sample size here once the collection and matching are complete.]

Our second empirical contribution establishes that the evaluation task is being actively reorganized rather than passively inheriting a change in volume. Comparing pull requests merged into the same repository in the same month, machine-authored changes attract more distinct human reviewers and are substantially less likely to be merged with no human review at all than comparable human-authored changes. Dependency-update bots in the same repositories show the opposite pattern, which demonstrates that the design detects perfunctory review where it exists. The effect is concentrated in smaller projects, where developers otherwise routinely merge their own work without review and where formal review had the least to do; larger projects that already require review on every change show little difference. That heterogeneity is the contingency \citet{barley1986technology} would predict, and it discriminates between accounts: machine authorship appears to remove an author's accountability that had partially substituted for formal review, rather than to change review capacity.

Three constructs must be kept apart throughout, and conflating them has caused difficulty in adjacent work. Aggregate review volume measures how much recorded activity a project has. A reviewer's share measures how that activity is distributed across people. Absolute counts measure what any individual actually did. A share can rise because one reviewer does more, because others do less, or because fewer participants are observed, and these have different organizational meanings. None of the three is a measure of burden: effort, difficulty, interruption, and available capacity are not recoverable from counts, and a concentration statistic cannot establish overload or strain \citep{tarafdar2007technostress}. We use concentration to describe allocation, and treat burden only as a limitation.

The stake, following \citet{zuboff1988smartmachine} and \citet{beane2019shadow}, is the distribution of learning rather than the distribution of effort. If evaluation concentrates on incumbents, the participants who would otherwise have reviewed lose the exposure through which code-base knowledge is acquired, and the cost appears later as a succession problem rather than immediately as an overload problem. We measure the allocation and the participation that bear on this; we do not measure learning, and we do not claim it.

Our research questions are:
\begin{enumerate}
\item \textbf{RQ1:} Does the human review work of a project concentrate on fewer reviewers after first detected AI adoption, once the estimate is identified against mean reversion in a bounded outcome?
\item \textbf{RQ2:} Does machine-authored code receive the same human evaluation as comparable human-authored code within the same project and period?
\item \textbf{RQ3:} Do concentration and differential evaluation coincide with changes in absolute activity, participation, and the continuity of pre-adoption reviewer identities?
\end{enumerate}

The paper contributes in three ways. It brings a model of knowledge-based hierarchies to bear on a complementary task that AI makes more necessary without performing, and tests its prediction on a process observable at the level of individual acts. It documents a specific and repairable identification failure that affects difference-in-differences estimates of bounded reverting outcomes, and supplies the placebo diagnostic that detects it. And it shows that evaluation of machine-authored work is actively differentiated by projects rather than waved through, which is the opposite of the prevailing expectation and bears directly on how review policy should respond.

\section{Related Literature and Conceptual Framework}
\label{sec:framework}
\subsection{AI, coordination, and workload}

The closest AI study is \citet{song2026impact}. Its review-activity results make clear that review provision is already part of the literature. Our proposed extension pairs within-repository reviewer shares with persistence of pre-adoption reviewer identities. It does not claim that earlier work ignored reviewers or organizational structure. The enterprise studies above address review capacity and production bottlenecks under different treatment definitions and institutional settings. Their findings motivate measuring downstream adaptation, but do not imply a common sign for review latency across all settings.

A related line of work examines changes in task allocation. \citet{hoffmann2025generative} study how Copilot access changes developers' allocation between coding, management, and collaborative activities. That question concerns the mix of activities undertaken by workers. Our question concerns the distribution of one activity across workers and the continuity of its leading participants. These are complementary dimensions: a task can become less prominent in an individual's work while its provision becomes more concentrated within a project.
The broader technology-and-work literature provides the framing for this distinction. Firm-level evidence links information technology to workplace reorganization and the deployment of skilled labor rather than to a single displacement channel \citep{bresnahan2002information}. The task-based view likewise distinguishes activities that technology displaces from tasks that it creates or reinstates \citep{acemoglu2019automation}. We bring that logic to a repeated project-level coordination task: the relevant outcome is who performs review and how stable that allocation is after adoption.

\subsection{Expertise coordination and the allocation of review}

\citet{faraj2000expertise} distinguish possessing expertise from coordinating its use. \citet{kudaravalli2017configural} further show why centralized and decentralized collaboration should not be evaluated through a single universal ordering. These perspectives motivate treating review allocation as an organizational outcome without assuming that more concentration is necessarily harmful.

The software-review literature likewise already addresses distribution. \citet{hajari2023factoring} examine unequal review activity and tradeoffs among reviewer expertise, workload distribution, and knowledge retention. Their recommendation setting differs from our adoption setting, but it directly rules out a novelty claim based simply on measuring review inequality. Our question concerns how an observed technology transition relates to the distribution of review activity and the persistence of pre-existing leaders. We do not equate a high review count with expertise, a monthly leader with a maintainer, or leadership turnover with knowledge loss.

\subsection{Four distinct constructs}

Let $x_{rit}$ be the number of distinct merged pull requests reviewed by human reviewer $r$ in repository $i$ during month $t$. A reviewer contributes at most one incidence to each pull request, regardless of repeated submissions. Let $V_{it}=\sum_r x_{rit}$ and $s_{rit}=x_{rit}/V_{it}$ when $V_{it}>0$. The monthly top-two share is
\begin{equation}
 C^{(2)}_{it}=s_{(1)it}+s_{(2)it},
\end{equation}
where parentheses denote ranks determined anew each month. This describes concentration in the observed activity, not a fixed pair of people.

Let $L_i$ be the person with the highest aggregate incidence count in the pre-adoption window. Incumbent persistence is
\begin{equation}
 P_{it}=\mathbf{1}\{\operatorname{leader}_{it}=L_i\}.
\end{equation}
A decline in $P_{it}$ means that the incumbent less often occupies the leading position.

Absolute activity, $x_{rit}$, is a third construct. For example, a reviewer who continues to review 10 pull requests while the repository total falls from 20 to 15 has an increasing share with no increase in counted activity. Conversely, doubling every reviewer's count leaves all shares unchanged. Burden is a fourth construct: its assessment would require effort, difficulty, time pressure, or capacity measures beyond $x_{rit}$. Process latency is elapsed calendar time and is not a substitute for active work time.

\subsection{Competing explanations}

AI adoption could reinforce incumbent reviewers if project-specific knowledge becomes more valuable for evaluating unfamiliar contributions. It could instead be accompanied by changes in task mix, reviewer availability, or entry that shift the busiest role to another participant. Automation could reduce the demand for some human reviews, while expanded production could increase it elsewhere. These possibilities do not provide a unique directional prediction for concentration, incumbent persistence, or latency without additional assumptions.

The three questions stated in the Introduction organize the analysis around concentration, the differential evaluation of machine-authored work, and absolute activity and continuity. Earlier project documents used different hypothesis numbering, and the analysis evolved through exploratory follow-ups; the current structure should not be interpreted as evidence of ex ante registration.

\section{Data and Measurement}
\label{sec:data}
\subsection{Sample and detected adoption}

The candidate pool spans Python, JavaScript, TypeScript, Go, Java, and Rust, stratified into four creation-year cohorts from 2018 through 2025. Discovery applies repository activity and quality filters, including a minimum of 10 stars and exclusion of forks. There are 8,972 candidates with the covariates required for matching. The analysis uses exact language and creation-cohort strata, then nearest-neighbor matching without replacement on standardized $\log(1+\mathrm{stars})$ and $\log(1+\mathrm{contributors})$. The main distance caliper is 0.25, yielding 2,320 pairs. The pipeline identifies 1,771 pairs with cache entries for both repositories. This eligibility rule does not guarantee usable observations at every horizon or for every outcome.

Treatment timing is the earliest detected signal among agent-associated PR authorship or branch names, AI-associated commit trailers or footers, and configuration-file signals, including references in \texttt{.gitignore}. These signal types provide different strengths of evidence. Configuration intent is not equivalent to observed use, and first detection need not be the actual adoption date. Controls are repositories without detected adoption, not confirmed non-users. Tool-specific first-use dates were not retained separately, limiting interpretation of overlapping tool categories.

The matching covariates are repository attributes observed during collection rather than a verified historical snapshot before treatment. Good balance on those attributes does not resolve selection into adoption and may involve post-treatment variables. A causal analysis should reconstruct pre-treatment covariates and reassess matching. The current draft reports estimates from the existing design without claiming that this concern has already been addressed.

\subsection{Review shares, identities, and coverage}

The original repository-month pipeline caps extraction at 40 merged PRs per repository-month. The subsequent PR export contains 293,095 non-sentinel PR rows across 11,757 pair-role-month cells, representing 1,405 pair IDs and 2,092 repositories. Applying a hypothetical 40-row cap to this export would remove 107,243 rows (36.6\%); 1,922 cells (16.4\%) exceed 40 rows. These are export-level counts, not an estimate of the historical cap's causal effect or a count of lost review incidences. The export's largest cell contains 1,000 PRs, and removal of the 40-PR cap does not prove exhaustive upstream retrieval. Further restrictions produce the 300-pair concentration panel and the 86-pair activity panel.

Reviewer-identity extraction also counts each human reviewer once per PR. It permits identifying an incumbent using months $h=-6,\ldots,-2$, then comparing that identity with the monthly leader. The shared-reviewer diagnostic has 309 usable matched pairs, of which 176 have no observed reviewer overlap between the control repository and any treated repository in the available identity data. Other identity analyses have different eligibility rules. Their effective sample sizes must be exported separately before final reporting.

The PR-level \path{reviewer_event_counts} field, however, counts individual review submissions. Existing Gini, rank-by-rank, and cumulative analyses built from that field therefore concern event-weighted activity. They are not automatically replications of incidence-based shares. Gini is defined over a specified roster and weighting scheme:
\begin{equation}
 G_{it}=\frac{\sum_{r=1}^{n}\sum_{q=1}^{n}|x_{rit}-x_{qit}|}{2n\sum_{r=1}^{n}x_{rit}}.
\end{equation}
For a fixed finite roster, its maximum is $(n-1)/n$, not exactly one. Including only active reviewers differs substantively from including a fixed pre-adoption roster with later zeros. Neither choice measures team membership without additional information.

\subsection{Process measures and interpretation}

The stored first-review outcome uses the earliest submitted review timestamp regardless of reviewer type. We therefore label it \emph{time to first recorded review}, not time to first human review. The variable named \path{review_rounds} counts submitted review records, including bots. It does not reconstruct human review--response cycles or distinguish substantive feedback from approvals. We label it \emph{review submissions per PR}. Time to merge is elapsed time from creation to closure for merged PRs. These definitions follow the current extraction code and do not change the stored estimates.

The dataset conditions on merged pull requests. Open and closed-unmerged PRs are outside the analysis, so changes in merge selection can affect the measured process. Repository-month means weight observed months equally, rather than weighting all PRs equally across the sample. Reviewers per PR measures participation on individual PRs. A null estimate for that outcome does not establish constant reviewer-pool size, constant membership, or equal oversight quality.

\section{Empirical Strategy and Inferential Status}
\label{sec:strategy}

The study uses the imputation approach of \citet{borusyak2024revisiting}. We interpret the estimates as associations given unresolved adoption selection and measurement concerns. A repository and calendar-month fixed-effects model is fit to control observations and treated observations before detected adoption:
\begin{equation}
 Y_{it}(0)=\alpha_i+\lambda_t+\varepsilon_{it}.
\end{equation}
Predicted untreated outcomes are subtracted from treated post-adoption observations. The pooled estimate averages these residuals over usable treated repository-months, and event-time estimates average within each horizon. The intended window is $h=-6,\ldots,-2$ before adoption and $h=0,\ldots,11$ after adoption. Month $h=-1$ is excluded, and September 2026 is excluded as incomplete at collection. Matched controls receive their treated partner's event date.

A causal interpretation requires appropriate untreated trends, no anticipatory outcome changes, adequate support for imputation, and assumptions about interference and adoption selection. Neither matching nor an insignificant pre-period diagnostic establishes these conditions. Detection errors, shared reviewers, evolving project activity, and contemporaneous governance changes are plausible threats.

There are also implementation limits. The joint pre-trend diagnostic sums squared lead-specific $t$ statistics and compares the sum with a chi-squared distribution. It omits cross-lead covariance, and its pre-period residuals come from observations used to fit the model. The statistic is not a fully implemented covariance-aware BJS pre-trend test. The supplementary analyses instead estimate lead indicators in an untreated-only fixed-effects regression and use the full matched-pair-cluster covariance matrix for a joint test. This diagnostic assesses differential pre-period changes without establishing adoption exogeneity.

The analytic ATT standard error is also an approximation: it uses the dispersion of repository-average residuals, whereas the pooled ATT averages repository-month residuals. The completed 1,000-replicate matched-pair bootstrap refits the imputation model and provides a more relevant uncertainty assessment for five pooled outcomes. Matching is held fixed. This does not address dependence across pairs connected by reviewers or validate uncertainty for every follow-up outcome.

Table~\ref{tab:main} uses the completed 1,000-replicate bootstrap. Its p-values are normal-approximation calculations from the reported ATT and bootstrap standard error, rather than bootstrap-tail probabilities, and are therefore approximate. Identity and subgroup results retain their separately documented inference. We avoid significance stars across incompatible inference procedures and treat the many exploratory comparisons as exploratory ; the new activity analysis separately reports Holm adjustment across its 14 outcome--specification tests.

\section{Results}
\label{sec:results}
\subsection{RQ1: Monthly concentration and its sensitivity}

\begin{table}[htbp]
\centering\small
\caption{Historical capped-panel estimates: contextual comparisons}
\label{tab:main}
\begin{tabularx}{\textwidth}{Xrrr}
\toprule
Outcome & ATT & Bootstrap SE & Approx. $p$\\
\midrule
Top-one human review share & 0.0173 & 0.0087 & 0.047\\
Top-two human review share & 0.0151 & 0.0055 & 0.006\\
First recorded review (hours) & $-65.0128$ & 31.1447 & 0.037\\
Merge time (hours) & $-62.8121$ & 27.5755 & 0.023\\
Review submissions per PR & 0.1804 & 0.0821 & 0.028\\
\bottomrule
\end{tabularx}
\par\smallskip{\footnotesize\noindent Source: Authors' calculations from the original capped extraction. Share outcomes count distinct human reviewer--PR incidences. Process outcomes follow the definitions in Section~\ref{sec:data}. Outcome-specific observation counts differ. \par}
\end{table}

The principal reconstructed concentration estimate is +0.0243 for top-two share (pair-bootstrap SE 0.0087; 95\% interval [0.0072, 0.0415]), using 300 pairs, 4,924 repository-month observations, and 1,571 treated post-adoption observations. HHI and active-reviewer Gini are imprecise on these same observations. A broad bot-name exclusion sensitivity retains 280 pairs and gives +0.0219 for top-two share (95\% interval [0.0032, 0.0406]); the filter is heuristic rather than validated account classification. Appendix~\ref{app:followup} describes the construction. Table~\ref{tab:main} preserves the original capped-panel estimates as historical comparisons, not as interchangeable estimates from the reconstructed sample. Tail share is the complement of top-two share and adds no independent test.

On the same 802-treated-repository subsample, removing the sampling cap changes top-two share from 0.0153 to 0.0082 and top-one share from 0.0141 to 0.0070. The uncapped top-two estimate has an analytic $t$ statistic of 2.47, but the existing notes also flag pre-period irregularities. These results show measurement sensitivity; the uncapped estimate is not a final benchmark.

The common-horizon diagnostic likewise matters. Restricting to $h=0,\ldots,5$ changes the pooled top-two estimate from approximately 0.0151 to 0.0081. At longer horizons, fewer cohorts contribute. The full-window estimate therefore combines dynamics with changing cohort composition.

\subsection{RQ2: Differential evaluation of machine-authored work}

This analysis compares pull requests merged into the same repository in the same
month, so that every comparison holds the project and the period fixed. Author
type is classified as an AI coding agent, a dependency-update bot, or a human.
Standard errors are clustered by repository. Two falsification checks accompany
the estimates: a random split of human-authored pull requests, for which the
true difference is zero by construction, and the dependency-update bots, whose
pull requests are widely understood to receive perfunctory review and which
therefore calibrate what low scrutiny looks like in these same repositories.

[RESULTS PENDING: report the agent-versus-human differences in distinct human
reviewers, the probability of merging with no human review, review events per
unit of change, and elapsed time to merge, each with and without adjustment for
change size; then the dependency-bot contrast, the human-split placebo, and the
moderation by project size. Estimate on the re-collected panel so that this
section and RQ1 rest on one consistently measured dataset.]

\subsection{RQ3: Absolute activity, participation, and incumbent continuity}

Earlier identity analyses reported reduced persistence, including an estimate near $-0.077$ with analytic SE 0.0151. These estimates do not use the independent historical selection design and are retained only as exploratory context. Our interpretation relies on the separated-window check below and the new absolute-activity decomposition, rather than treating the older estimate as established displacement.

Selecting the incumbent as the pre-period maximum raises measured pre-period persistence by construction. Both treated and control repositories show a decline around the event boundary. Symmetric construction can remove a common component, but a similar visual decline does not establish equal selection bias. An independent-window analysis selects the leader in months $h=-6,-5,-4$ and uses $h=-3,-2$ as the estimation baseline. On 135 eligible pairs, persistence falls by 0.0725 (bootstrap SE 0.0432; 95\% interval [$-0.1571$, 0.0121]). The negative direction remains, but this check does not establish a statistically precise displacement effect.

The evidence does not establish what arrangement follows displacement. The full-sample month-to-month leader-change estimate is negative, but it becomes approximately zero in the 176-pair restriction ($t=-0.02$). Pre-adoption top-two overlap also becomes statistically indistinguishable from zero there. Moreover, post-period top-two identity overlap across two six-month blocks does not differ detectably between treated and control repositories in the descriptive comparison. The data therefore do not identify a stable handover or increased rotation.

The incumbent's absolute activity offers a direct check on interpretation. In the 309-pair displacement follow-up, the incumbent's review count has an ATT of $-0.265$ with bootstrap $p=0.703$. The incumbent-share estimate is also imprecise and fails the existing pre-period diagnostic. Descriptively, 43.6\% of replacement-leader observations in treated repositories involve the former second-ranked reviewer, compared with 31.5\% in controls. These are selected replacement months, not a causal comparison of successor types, and rank does not identify tenure or knowledge. The combination illustrates why a change in who ranks first need not mean that the former lead stops working.


\subsubsection{New activity decomposition.}
We freeze the historical top two and the observed reviewer roster using $h=-6,-5,-4$, exclude cores with fewer than two reviewers or a tie across the second-place boundary, and estimate using $h=-3,-2$ and $h=0,\ldots,11$. Both pair members must have usable pre- and post-adoption information. This leaves 86 pairs and 1,260 repository-month observations, including 516 treated post-adoption observations. Counts use one incidence per reviewer per merged PR. ``Outside roster'' means absent from this short historical window, not necessarily new to the project. The design remains conditional on observed merged-PR months; missing months cannot be treated as verified zeros.

\begin{table}[htbp]
\centering\small
\caption{Absolute activity and participation: 86-pair historical-core sample}
\label{tab:activity}
\begin{tabularx}{\textwidth}{Xrrrr}
\toprule
Outcome & Baseline & Estimate & 95\% interval & Holm $p$\\
\midrule
Historical top-two incidences & 18.652 & 0.742 & [-2.927, 4.410] & 1.000\\
Other historical-reviewer incidences & 17.037 & 0.979 & [-2.265, 4.224] & 1.000\\
Outside-roster incidences & 3.177 & 2.032 & [0.465, 3.600] & 0.155\\
Total reviewer--PR incidences & 38.866 & 3.753 & [-2.791, 10.298] & 1.000\\
Active reviewers & 5.683 & 0.281 & [-0.467, 1.029] & 1.000\\
Fraction of PRs reviewed & 0.756 & -0.009 & [-0.069, 0.050] & 1.000\\
Merged PRs & 29.293 & 4.668 & [-0.114, 9.450] & 0.724\\
\bottomrule
\end{tabularx}
\par\smallskip{\footnotesize\noindent Notes: Baseline means use treated pre-period observations. Normal intervals use standard deviations from 1,000 matched-pair bootstrap refits (seed 20261102; matching fixed). Holm adjustment covers seven outcomes in each of two specifications, including the historical-roster-of-three-or-more restriction. Outside-roster activity fails its pre-period diagnostic ($p=0.0136$). These intervals do not account for dependence across pairs.\par}
\end{table}

Historical-core activity is +0.74 incidences per month and other historical-reviewer activity is +0.98, with wide intervals spanning decreases and increases. Active reviewers and the fraction of PRs reviewed also have imprecise estimates. Merged-PR volume is +4.67 (95\% interval [$-0.11$, 9.45]), so an expansion in volume is possible but not established. The only nominally significant endpoint, outside-roster incidences (+2.03; unadjusted $p=0.0110$), fails the pre-period diagnostic and has Holm $p=0.1547$. No endpoint rejects after adjustment across the 14 comparisons.

Restricting to historical rosters of at least three reviewers leaves 43 pairs and 680 observations. No endpoint rejects after adjustment, and outside-roster activity again fails its pre-period diagnostic ($p=0.0100$). These results do not establish either a shift toward the historical core or a contraction in peripheral participation. Their wide intervals also do not establish the absence of such changes. Because the 86-pair sample is selected differently from the 300-pair concentration panel, the decomposition cannot be treated as a direct explanation of that panel's estimate. Appendix~\ref{app:activity} reports all subgroup endpoints.

\subsection{Process outcomes}

Table~\ref{tab:main} shows shorter elapsed time to a first recorded review and to merge, alongside more review submissions. These are calendar and recorded-review measures; they include automation, queues, and author response time.

Existing PR-size and volume adjustments leave the latency point estimates broadly similar in their respective subsamples, although significance varies by specification. Adjustment for post-adoption PR characteristics changes the estimand and is not a definitive mechanism test. On the 296-repository, complete-pre-history subsample, which is busier on average, the signs reverse: merge time increases by about 40.7 hours and review submissions decrease by about 0.146. The restriction selects a distinctive group. Without a direct, comparably estimated moderator test, this should be described as sample sensitivity rather than proof that volume causes the reversal.

\subsection{Reviewer overlap and heterogeneous estimates}

Of the 309 pairs in the reviewer-overlap diagnostic, 133 (43\%) have a control reviewer who also appears in a treated repository. This establishes observed cross-arm connectivity and a possible interference channel, not proof that spillovers occurred. The 176-pair restriction removes observed overlap within the available identity data, not all possible exposure.

\begin{table}[htbp]
\centering\small
\caption{Selected estimates in the 176-pair sample without observed reviewer overlap}
\label{tab:clean}
\begin{tabularx}{\textwidth}{Xrr}
\toprule
Outcome & ATT & Analytic $t$\\
\midrule
Top-two human review share & 0.0208 & 2.79\\
Incumbent remains monthly leader & $-0.1106$ & $-3.70$\\
Monthly leader changes & $-0.0006$ & $-0.02$\\
First recorded review (hours) & $-89.7650$ & $-2.92$\\
Merge time (hours) & $-56.1510$ & $-1.88$\\
Review submissions per PR & 0.1756 & 1.82\\
\bottomrule
\end{tabularx}
\par\smallskip{\footnotesize\noindent Source: Authors' calculations. Analytic inference is not directly comparable with Table~\ref{tab:main}'s bootstrap inference. The restriction also changes the population of repositories.\par}
\end{table}

Top-two share remains positive in this restriction. However, relative to all 309 identity-covered pairs, its point estimate falls from 0.0318 to 0.0208. Comparing the restricted estimate only with the full 1,771-pair estimate would obscure that difference. The restriction cannot distinguish interference, heterogeneity, and selection.

Adoption-mode comparisons suggest a larger top-two-share change among human-assisted than agent-associated adopters in one bootstrap comparison ($p=0.037$). Tool-specific and adoption-era differences are generally imprecise, while several within-group estimates change sign across matching calipers. Tool buckets overlap and use a combined first-detection date. We therefore do not infer equivalent effects across tools or generations. Likewise, subgroup significance in one group and insignificance in another does not establish a difference between them.

\subsection{Exploratory distributional and mechanism checks}

Event-weighted cumulative analyses do not detect an increase in the top-two share over the entire observation window, although a cumulative Gini estimate is positive. The direct monthly--cumulative gap comparison is itself not significant. Consequently, a significant monthly estimate and an insignificant cumulative estimate do not establish increased rotation or a difference between the two effects.

The rank-by-rank event-weighted analysis reports declining shares at ranks three and four. It conditions each rank on that rank being present, producing different samples across ranks. Adding the rank-one and rank-two ATTs therefore does not recover an ATT for top-two share. Differences from the main incidence-based estimate cannot be attributed solely to the cap. A common-roster, common-sample incidence analysis is needed before claiming that the middle of the reviewer distribution is being displaced.

Other exploratory paths do not currently support mechanism claims. Episodic incumbent inactivity fails matching and bootstrap sensitivity. Permanent exit and incumbent authoring changes are not detected. Entrant-share and reciprocity constructions fail existing pre-period diagnostics. Leadership-tenure differences are entangled with pre-period concentration. These results limit interpretation but do not prove that exit, reallocation, entry, or network changes are absent.

\section{Discussion}
\label{sec:discussion}

The principal finding is a positive association between detected adoption and the monthly top-two share of human review incidences. The new contribution is the direct examination of absolute activity alongside that relative measure. Neither the activity decomposition nor independent-window persistence establishes who absorbs or relinquishes review work. A smaller denominator, changes in identity, changing merged-PR composition, or shifts among active reviewers remain possible. The analysis therefore distinguishes an observed concentration pattern from an identified organizational mechanism.

The conceptual implication is that relative allocation, absolute activity, identity continuity, and burden require separate evidence. Our results illustrate the empirical importance of that distinction rather than demonstrate a particular reorganization of work. For managers, concentration may motivate investigation of capacity and knowledge coverage, but the present estimates do not justify a staffing, workload, or quality diagnosis.

\section{Limitations and Next Tests}

Several limitations are central to the interpretation. First, visible adoption signals are incomplete and heterogeneous, and collection-time matching covariates may be affected by adoption. Second, merged-only data condition on a potentially changing selection process. Third, capped sampling, event versus incidence weighting, and varying coverage make some apparently comparable estimates different estimands. Fourth, incumbent selection and tied ranks can mechanically affect persistence. Fifth, shared reviewers challenge independent-pair inference. Sixth, approximate pre-period tests and numerous exploratory comparisons preclude treating nominal significance as definitive causal validation.

The reconstruction and activity decomposition complete part of the measurement agenda, but several checks remain open. The estimator has not yet been validated against an external maintained implementation. Extraction completeness, zero-review versus missing-month classification, and validated account types require further audit. Historical matching variables, independently measured adoption dates, longer identity-selection windows, and inference allowing dependence across pairs would strengthen the design. Fixed follow-up cohorts should be defined by verified coverage rather than positive reviewing. The new activity estimates use the existing observed-month weighting over months 0--11, not the proposed equal-repository primary estimand over months 0--5.

A mechanism extension should compare pre-adoption expertise and availability across the full reviewer candidate pool, then examine who becomes a successor. A burden extension would need individual review counts, task complexity, active effort or time-pressure measures, and capacity information. Linking allocation to reverts, defects, retention, or knowledge coverage would establish whether the observed restructuring has beneficial or harmful consequences. These are proposed extensions, not outcomes established in the current dataset.

\section{Conclusion}

First detected AI adoption is associated with higher top-two review share in the reconstructed 300-pair sample. The new absolute-activity analysis does not establish whether that pattern reflects more work by historical core reviewers, less work by other reviewers, or reduced participation. No activity endpoint survives multiplicity adjustment, and the nominal outside-roster increase fails a pre-period diagnostic. Independent-window incumbent persistence remains uncertain. The evidence supports a focused concentration association, not a demonstrated displacement, burden, or productivity mechanism. Its causal magnitude and organizational consequences remain open.

\bibliographystyle{informs2014}
\bibliography{references}

\begin{APPENDIX}{Supplementary Empirical Analyses}
\label{app:followup}
\renewcommand{\theHsection}{supplement.\arabic{section}}

\section{Common-Sample Reconstruction and Inference}
We reconstruct monthly reviewer counts from distinct reviewer identities on each merged PR. Repeated submissions by the same reviewer on a PR contribute one incidence. We exclude two empty extraction sentinel rows, validate repository identities and event dates against the matched pairs, exclude September 2026, and retain $h=-6,\ldots,-2$ and $h=0,\ldots,11$. Each retained pair must have an observed review month before and after adoption in both arms. Months without recorded reviews have undefined shares and are not assigned zero. The resulting concentration panel contains 300 pairs, 4,924 repository-month observations, and 1,571 treated post-adoption observations.

All four concentration measures use exactly this panel. HHI is $\sum_r s_{rit}^2$; Gini is computed over active reviewers. A single-reviewer month has Gini zero and HHI one, reflecting the measures' different treatment of participation. The reconstruction removes the monthly 40-PR cap but cannot establish exhaustive upstream retrieval.

We refit the repository and calendar-month effects in each of 1000 matched-pair bootstrap draws, holding matching fixed. Reported intervals are normal approximations using bootstrap standard deviations. The pre-period diagnostic estimates treated lead indicators using untreated observations, with the earliest observed pre-period month omitted, and tests them jointly using the full pair-cluster covariance matrix. These procedures do not account for reviewer dependence across pairs.

\section{Task Allocation and Reviewer Relationships}
Studies of technology and organization distinguish task assignment from aggregate productivity \citep{bloom2014distinct,hoffmann2025generative}. We therefore compare reviewer HHI with HHI over merged-PR authors in the same repository-month and directly estimate their difference. We also measure the share of review incidences supplied by individuals who author at least one merged PR that month. This is overlap between observed roles, not time allocation. Finally, we measure the share of author--reviewer incidences belonging to a pair observed together on at least two PRs in the month, excluding self-review links. Repeated links measure relationship reuse rather than review quality.

Requiring valid author and reviewer distributions and at least one non-self author--reviewer link gives a common panel of 290 pairs and 4,701 months. We repeat this analysis on the 24 pairs observed in all six post-adoption months $h=0,\ldots,5$. A separate account-name sensitivity excludes names containing bot, assistant, copilot, coderabbit, gemini, devin, dependabot, or renovate. This deliberately broad filter is a sensitivity check, not validated account classification; it leaves 280 pairs. Top-two share remains positive under that filter (0.0219, SE 0.0095), whereas Gini is imprecise (0.0196, SE 0.0118). The six-month top-two estimate is 0.0400 (SE 0.0224); its small selected sample limits precision.

\section{Independent Historical Selection}
Research on ownership and expertise motivates distinguishing continuity from distribution \citep{bird2011ownership}; evidence of learning from AI also motivates examining established participant relationships \citep{brynjolfsson2025work}. We identify the incumbent, top-two reviewers, reviewer roster, and author--reviewer links using only observations from $h=-6,-5,-4$. We then estimate outcomes using $h=-3,-2$ and $h\geq0$, leaving historical selection observations out of estimation. Tied historical leaders are excluded. If several reviewers tie for first in a subsequent month, an incumbent among them receives persistence weight equal to the reciprocal of the number tied.

The 135-pair panel contains 1,924 months and 736 treated post-adoption observations. Besides incumbent persistence, we estimate the share supplied by the historical top two, the share supplied by reviewers outside the historical roster, and the share of links belonging to historical author--reviewer pairs. The short historical window measures observed prior participation rather than expertise or true newcomer status. The two-month estimation baseline also limits the power of pre-period testing.

\begin{table}[htbp]
\centering\small
\caption{Selected supplementary estimates}
\label{tab:followup}
\begin{tabularx}{\textwidth}{Xrrr}
\toprule
Outcome & Pairs & Estimate & Bootstrap SE\\
\midrule
Top-two share, reconstructed incidence panel & 300 & 0.0243 & 0.0087\\
Gini, same incidence panel & 300 & 0.0205 & 0.0110\\
HHI, same incidence panel & 300 & 0.0098 & 0.0154\\
Review HHI minus author HHI & 290 & 0.0104 & 0.0186\\
Share supplied by concurrent PR authors & 290 & 0.0259 & 0.0233\\
Repeated author--reviewer link share & 290 & 0.0199 & 0.0258\\
Incumbent persistence, independent history & 135 & $-0.0725$ & 0.0432\\
Historical top-two share & 135 & 0.0045 & 0.0331\\
Outside historical roster share & 135 & $-0.0018$ & 0.0268\\
Historical author--reviewer link share & 135 & $-0.0580$ & 0.0370\\
\bottomrule
\end{tabularx}
\par\smallskip{\footnotesize\noindent Notes: Estimates are in outcome units; multiplying share outcomes by 100 gives percentage points. Standard errors use 1000 pair-bootstrap refits. Task-allocation and historical outcomes are exploratory. Holm adjustment across their 16 outcome--specification tests yields no rejection at 5\%.\par}
\end{table}

These analyses support the top-two concentration pattern but do not identify a change in the separation of coding and reviewing, relationship reuse, or participation by previously unobserved reviewers. The independent-window persistence estimate retains its negative direction with substantial uncertainty. Neither insignificant coefficients nor historical participation measures establish the absence of reorganization or the loss of expertise.

\section{Coverage and Dependence Diagnostics}
The retained-versus-excluded diagnostic identifies 300 eligible pairs among 1,396 with non-September PR observations. Retained pairs have more collection-time contributors (82.4 versus 48.9; standardized difference 0.627) and more observed baseline PR rows across both arms (152.3 versus 72.9; standardized difference 0.538). This is evidence of selective coverage, not proof that activity is measured exhaustively or that findings generalize only to large repositories.

A broader baseline network diagnostic finds 448 reviewer identities shared across pairs, touching 408 of 1,136 pairs (35.9\%); the largest connected component contains 209 pairs. This network is broader than either estimation sample. Shared identities create a possible channel for correlated errors or interference, but the connectivity counts do not estimate error covariance or establish the direction or magnitude of a standard-error correction. We therefore treat independent-pair inference as a limitation, rather than assert that every reported interval is necessarily too narrow.

The September-to-August observed PR-volume ratio has a median of 0.486 among 635 repositories with at least five August rows; 204 have no September rows in the export. These descriptive counts are consistent with incomplete coverage but cannot distinguish truncation from activity changes or event-window eligibility. September remains excluded under the existing extraction rule; this ratio alone does not verify a collection cutoff.

The new export comparison joining on pair ID, role, year, and month was not valid: 2,277 of its 2,494 nominal matches refer to different repositories and 2,278 have different event dates. Its reported capped--uncapped differences are excluded from the evidence. A valid comparison must reconcile repository identity, event date, and extraction definitions before attributing differences to the cap. Likewise, the cache-based baseline-balance diagnostic aggregates by repository rather than matched event date and uses available observations separately in each arm; it is not a validated pair-complete balance table for the reconstructed sample. Finding all retained matches within the chosen caliper does not establish that the caliper was non-binding during matching.

For the 86-pair activity sample over months 0--5, equal-repository and observed-month point estimates are close for core incidences (2.285 versus 2.278) and total incidences (7.017 versus 6.825). This is a limited point-estimate sensitivity, not a new inference result or proof of invariance for all endpoints: the reviewed-PR fraction changes sign ($-0.0006$ versus 0.0097). Neither weighting comparison resolves the difference between the 0--5 and 0--11 horizons.

\section{Activity Subgroup Results}
\label{app:activity}
\begin{table}[htbp]
\centering\small
\caption{Historical roster of at least three reviewers: 43 pairs}
\begin{tabularx}{\textwidth}{Xrrrr}
\toprule
Outcome & Estimate & 95\% interval & Pre-period $p$ & Holm $p$\\
\midrule
Historical top-two incidences & 1.848 & [-3.317, 7.012] & 0.573 & 1.000\\
Other historical-reviewer incidences & 1.711 & [-3.036, 6.458] & 0.673 & 1.000\\
Outside-roster incidences & 1.190 & [-0.841, 3.222] & 0.010 & 1.000\\
Total reviewer--PR incidences & 4.749 & [-4.381, 13.879] & 0.334 & 1.000\\
Active reviewers & -0.178 & [-1.419, 1.064] & 0.219 & 1.000\\
Fraction of PRs reviewed & 0.019 & [-0.050, 0.089] & 0.226 & 1.000\\
Merged PRs & 4.046 & [-3.081, 11.172] & 0.300 & 1.000\\
\bottomrule
\end{tabularx}
\par\smallskip{\footnotesize\noindent Notes: 680 observations, including 279 treated post-adoption observations. Definitions, bootstrap settings, and the 14-test Holm family match Table~\ref{tab:activity}. The short estimation baseline provides limited pre-period diagnostic power.\par}
\end{table}

\end{APPENDIX}
\end{document}
